\chapter*{ABSTRACT}
\addcontentsline{toc}{chapter}{ABSTRACT}
\noindent\rule{\linewidth}{2pt}

\vspace{0pt}

\vspace{30pt}

\noindent
The rapid growth of the aviation industry has increased the demand for efficient and reliable airline reservation systems. Traditional manual booking methods are often time-consuming, error-prone, and inefficient when handling large volumes of passenger data and flight schedules. To address these challenges, this project presents the design and development of an Airline Reservation System using Object-Oriented Programming (OOP) in Java.

\par
The proposed system provides a user-friendly platform that enables passengers to search for flights, check seat availability, book tickets, cancel reservations, and view booking history. The system is developed using Java for application logic and MySQL for database management, with JDBC used for seamless connectivity between the application and the database.

\par
The implementation of OOP concepts such as encapsulation, inheritance, abstraction, and polymorphism ensures modularity, scalability, and maintainability of the system. Encapsulation is used to protect sensitive data, inheritance supports code reusability, abstraction simplifies complex operations, and polymorphism enables flexible functionality.

\par
The system ensures real-time updates of seat availability and maintains accurate passenger and booking records. It reduces manual errors and improves the efficiency of reservation management. The results demonstrate that the proposed system provides a reliable and efficient solution for airline ticket booking and management.

\par
This project highlights the practical application of Object-Oriented Programming in solving real-world problems and lays the foundation for future enhancements such as online payment integration, web-based interfaces, and cloud deployment.